% Abstract (150--300 words) and keywords (3--5), in English.
\renewcommand{\abstractname}{Abstract}
\begin{abstract}
Local government staff in Spain review the Spanish Official Gazette (BOE) and the Official Gazette of Andalusia (BOJA) by hand every day to find the job calls, grants, deadlines and regulations that affect their municipality, a slow and error-prone task. This project presents a platform that lets them query those gazettes in natural language, with answers grounded only in the documents and citing every source.

The system applies retrieval-augmented generation (RAG). A Go service downloads and filters the gazettes every day; a .NET backend built on Semantic Kernel chunks them, indexes them in PostgreSQL and Qdrant, and generates the answers; and an Angular interface streams them in real time. The platform serves several municipalities with isolated data, includes an Arena to compare two versions of the assistant blindly, and records the usage, latency and cost of every model call. It is orchestrated with .NET Aspire and deployed with Docker Compose, continuous integration and OpenTelemetry-based observability.

As part of the project, four chunking strategies (fixed-size, recursive, semantic and hierarchical parent-child) are compared using local models, on collections of 17, 200 and 5,994 documents and with a question set grouped by document size and difficulty, evaluated with RAGAS and with judge-free retrieval metrics. The hierarchical strategy achieves the best retrieval in every scenario, with the largest advantage on long documents and on near-duplicate ones. The study also shows that evaluating only on the documents that answer the questions overestimates quality: the top-1 hit rate drops from 1.00 with 17 documents to between 0.55 and 0.73 with the full collection.

\medskip
\noindent\textbf{Keywords:} RAG, LLM, official gazettes, local government, chunking
\end{abstract}
